\section{总结与延伸}

本讲从三种反馈来源解释 Agent 如何改进。ReAct 在推理时交错 thought、action 和 observation，用工具反馈修正当前轨迹；RLEF 把代码执行结果作为可扩展 reward，通过多轮交互和 PPO 学习修复策略；Constitutional AI 把人类原则扩展为 AI critique、revision 和 preference labels，用于监督微调与 RLAIF。

最重要的统一结论是：\textbf{反馈的价值不取决于它是否由人或 AI 产生，而取决于它与真实目标的对应关系、错误粒度、可扩展性和抗投机能力。} 环境反馈通常更客观但领域受限，AI feedback 更通用但需要外部锚点。

下一讲讨论 planning 与 multi-step reasoning：当任务需要很长的行动链时，系统应如何搜索计划树、把串行计划转化为并行执行，以及如何用 synthetic trajectories 训练多步工具使用能力。

\end{document}
